% ===========================================================================
%  E11 — Site público e captação
% ===========================================================================
\chapter{E11 \textperiodcentered{} Site público e captação}
\label{ep:e11}

\begin{fichaepico}
\begin{campos}
  \item[Objetivo] Quem ainda não é cliente entende a proposta do Entre Redes e
    deixa contato; quem já é encontra a entrada e os documentos legais.
  \item[Para quem] Visitante (secretarias, escolas, docentes e famílias).
  \item[Valor] É a porta comercial do produto e o lugar onde a rede e a família
    leem o que a plataforma promete.
  \item[Estado] \parcial{} --- site e cadastro prontos; faltam os caminhos
    para o login e para os documentos legais.
\end{campos}
\end{fichaepico}

\begin{tabularx}{\linewidth}{@{}L{16mm}X L{20mm}@{}}
\toprule
\headrow \thd{ID} & \thd{Feature} & \thd{Estado}\\
\midrule
\id{F11.1} & Site ENTRE REDES & \parcial\\
\zebra \id{F11.2} & Cadastro de interesse & \pronto\\
\id{F11.3} & Política de Privacidade e Termos de Uso & \parcial\\
\bottomrule
\end{tabularx}

% ---------------------------------------------------------------------------
\feature{F11.1}{Site ENTRE REDES}{\parcial}

\begin{historia}{H11.1.1}{Conhecer a proposta}{\pronto}
\hu{Visitante de uma secretaria de educação}{entender o que é o Entre Redes e
como ele funciona}{decidir se vale conversar com a Unooria}
\aceite
\begin{itemize}
  \item O site tem nove seções: Programa, O desafio, Nossa resposta, O método,
    As missões, Gestão, Currículo, Docentes e Famílias.
  \item O botão \enquote{Explorar} abre o índice; os prints das missões podem
    ser ampliados.
  \item No celular, o convite encurta para \enquote{Cadastro}.
  \item Quem já está com sessão aberta e abre o site vai direto para a própria
    área.
\end{itemize}
\end{historia}

\begin{historia}{H11.1.2}{Achar a entrada e os documentos a partir do site}{\afazer}
\hu{professor da rede que chegou pelo site}{encontrar o botão de entrar e os
documentos legais}{não depender de saber o endereço da tela de entrada}
\aceiteproposta
\begin{itemize}
  \item O cabeçalho do site tem \enquote{Entrar}.
  \item O rodapé leva à Política de Privacidade e aos Termos de Uso.
\end{itemize}
\end{historia}

% ---------------------------------------------------------------------------
\feature{F11.2}{Cadastro de interesse}{\pronto}

\begin{historia}{H11.2.1}{Deixar o meu contato}{\parcial}
\hu{Visitante}{deixar os meus dados para a Unooria me procurar}{conhecer a
solução com calma}
\aceite
\begin{itemize}
  \item Os três botões \enquote{Quero conhecer a solução} abrem a mesma janela:
    Nome, E-mail profissional, Telefone / WhatsApp, \enquote{Rede, secretaria
    ou escola}, Seu cargo e a caixa \enquote{Autorizo o contato da Unooria
    sobre o ENTRE REDES}.
  \item Depois de \enquote{Enviar cadastro}: \enquote{Cadastro enviado. A
    equipe da Unooria entrará em contato.}
  \item A resposta é sempre a mesma --- inclusive para cadastro repetido ou
    barrado ---, para não revelar quem já se cadastrou.
  \item O texto exato da autorização é guardado junto com o cadastro.
  \item Há proteção contra robôs e limite de envios por hora.
\end{itemize}
\pendencia{a janela de cadastro não leva à Política de Privacidade.}
\end{historia}

% ---------------------------------------------------------------------------
\feature{F11.3}{Política de Privacidade e Termos de Uso}{\parcial}

\begin{historia}{H11.3.1}{Ler o que a plataforma faz com os dados}{\parcial}
\hu{família ou encarregado de dados da rede}{ler em linguagem clara o que a
plataforma coleta e por quê}{confiar na plataforma antes de a criança usar}
\aceite
\begin{itemize}
  \item A Política de Privacidade tem 17 seções e os Termos de Uso 16, com
    sumário lateral que acompanha a leitura.
  \item Os Termos têm \enquote{As experiências são simulações --- leia esta
    parte}, \enquote{O que cada perfil pode fazer} e \enquote{A inteligência
    artificial na plataforma}.
  \item Os dois documentos só afirmam o que a plataforma faz de verdade, e
    terminam com \enquote{O que ainda está sendo definido}.
  \item A Política de Privacidade está na versão 1.1, em vigor desde 25 de
    setembro de 2026. Os Termos de Uso, ainda um rascunho para revisão
    jurídica, foram atualizados no mesmo dia.
  \item Os dois já dizem as regras da A Seleção: o colega com quem a
    publicação foi pareada no Mural vê os três posts, a legenda e o primeiro
    nome do estudante; e a IA lê a legenda para sugerir ao professor um nível,
    que a criança não vê.
\end{itemize}
\pendencia{razão social e CNPJ da Unooria ainda em branco; e estão declarados
como indefinidos os prazos de guarda e eliminação, o registro de acesso
individual, o canal para responsáveis e o relatório de impacto. Os Termos ainda
citam o UnooTube, arquivado, como exemplo de aplicativo imitado: a frase está
com o PO para revisar.}
\end{historia}
